\documentclass[12pt]{article}
\usepackage[T1]{fontenc}
\usepackage[utf8]{inputenc}
\usepackage{lmodern}
\usepackage{textcomp}
\usepackage{enumitem}
\usepackage{geometry}
\geometry{margin=1in}
\begin{document}
\begin{center}
Answer Key - Health Sciences - K-12 Level Assessment 23 \
\small (Answers are ordered exactly as in the question paper.)
\end{center}

\section*{Multiple Choice}
\begin{enumerate}[label=\arabic*.]
\item \textbf{Answer:} D
\\ \textit{Options:} A) As simple Fe$_{2}$+ in the serum; B) Bound to albumin; C) Bound to ferritin; D) Bound to transferrin
\\ \textit{Note:} Iron is transported in the circulation from the intestine to the sites of metabolism in the body bound to transferrin, a specific protein that facilitates the safe and efficient delivery of iron to cells while preventing toxicity.
\item \textbf{Answer:} C
\\ \textit{Options:} A) \textless{}0.8 g protein per kg bodyweight per day; B) 0.8-1.2 g protein per kg bodyweight per day; C) 1.2-1.7 g protein per kg bodyweight per day; D) \textgreater{}2.0 g protein per kg bodyweight per day
\\ \textit{Note:} Athletes require 1.2-1.7 grams of protein per kilogram of bodyweight daily to support muscle repair, recovery, and overall performance, reflecting their increased protein needs compared to the general population.
\item \textbf{Answer:} C
\\ \textit{Options:} A) Docosahexaenoic acid can be synthesised from linolenic acid in vegans; B) Docosahexaenoic acid is absent from vegan diets; C) Lack of docosahexaenoic acid in vegans causes visual and cognitive impairment; D) Microalgae can be used to synthesise docosahexaenoic acid
\\ \textit{Note:} The statement is not true because while docosahexaenoic acid (DHA) is important for visual and cognitive function, vegans can obtain sufficient DHA through algae-based supplements or by converting alpha-linolenic acid (ALA) from plant sources, thus not necessarily causing impairment.
\item \textbf{Answer:} A
\\ \textit{Options:} A) Folate; B) Niacin; C) Riboflavin; D) Vitamin B$_{6}$
\\ \textit{Note:} Folate deficiency leads to megaloblastic anemia because this vitamin is essential for DNA synthesis and cell division, and its lack results in the production of larger, immature red blood cells.
\item \textbf{Answer:} A
\\ \textit{Options:} A) Carries a risk of colonic cancer; B) Can be treated with a LOFFLEX diet; C) Can be caused by small intestinal bacterial overgrowth; D) Can be caused by milk intolerance
\\ \textit{Note:} Ulcerative colitis is a chronic inflammatory bowel disease that increases the risk of developing colonic cancer due to prolonged inflammation and dysplasia in the colonic mucosa.
\end{enumerate}

\section*{True / False}
\begin{enumerate}[label=\arabic*.]
\setcounter{enumi}{5}
\item \textbf{Answer:} True
\\ \textit{Options:} A) True; B) False
\\ \textit{Note:} The statement is true because studies have shown that a significant percentage of individuals with alcohol use disorder develop alcoholic myopathy, with estimates indicating that between 40 to 60\% of alcoholics experience this muscle condition due to the toxic effects of alcohol on muscle tissue.
\item \textbf{Answer:} True
\\ \textit{Options:} A) True; B) False
\\ \textit{Note:} Intestinal failure can lead to significant fluid and electrolyte imbalances, and sodium supplementation of 90 mmol/l or more helps to maintain proper hydration and electrolyte levels, thereby supporting overall health and function.
\item \textbf{Answer:} True
\\ \textit{Options:} A) True; B) False
\\ \textit{Note:} Muscle glycogen is stored in muscles and provides a readily available source of glucose for quick energy during high-intensity activities such as sprinting, making it the primary fuel source for short-term muscle contractions.
\item \textbf{Answer:} False
\\ \textit{Options:} A) True; B) False
\\ \textit{Note:} Other diet-derived compounds, such as resveratrol, green tea catechins, and quercetin, also exhibit epigenetic effects, demonstrating that curcumin is not the only compound with such properties.
\item \textbf{Answer:} False
\\ \textit{Options:} A) True; B) False
\\ \textit{Note:} The statement is false because the worldwide prevalence of obesity is approximately 13\%, significantly lower than the claimed 39\%.
\end{enumerate}

\section*{Long Form Response}
\begin{enumerate}[label=\arabic*.]
\setcounter{enumi}{10}
\item \textbf{Answer:} Malnutrition in older adults significantly increases both mortality and morbidity, meaning it can lead to higher death rates and greater illness. When older adults do not get enough nutrients, their immune systems weaken, making them more susceptible to infections and diseases. For instance, a lack of protein can cause muscle loss, which affects mobility and increases the risk of falls, leading to serious injuries. Additionally, malnutrition can exacerbate existing health conditions, such as diabetes and heart disease, resulting in more frequent hospitalizations. Therefore, addressing malnutrition in older adults is crucial for improving their overall health and longevity.
\\ \textit{Note:} correct because it clearly articulates the direct link between malnutrition and increased mortality and morbidity in older adults, illustrating how inadequate nutrition weakens the immune system and exacerbates existing health issues, thereby providing relevant examples to support the analysis.
\item \textbf{Answer:} Non-esterified fatty acids are taken up by cells through two main mechanisms: diffusion and carrier-mediated transport. Diffusion allows these fatty acids to move freely across the cell membrane due to their lipid-soluble nature, which helps them enter cells without needing additional assistance. However, the uptake is often enhanced by carrier-mediated transport, which involves specific proteins that facilitate the movement of fatty acids into the cell. These transport proteins include fatty acid binding proteins and translocases, which help in efficiently transporting fatty acids across the membrane and into the cytoplasm. Together, these mechanisms ensure that cells can effectively absorb the fatty acids they need for energy and various metabolic processes.
\\ \textit{Note:} correct because it accurately describes the two primary mechanisms-diffusion and carrier-mediated transport-by which non-esterified fatty acids enter cells, emphasizing the role of their lipid-soluble nature and the involvement of specific transport proteins.
\end{enumerate}

\vfill
\noindent\textit{End of Answer Key}
\end{document}